% META: {"id": "1NSI-ALGO-PARCOURS-TRIS-REMED-06", "chapitre": "1NSI-ALGO-PARCOURS-TRIS", "type_objet": "remediation", "capacites": ["1NSI-ALGO-PARCOURS-TRIS-C6"], "status": "approved", "capacites_codes": ["C6"], "statut": "complete"}

%---------------------------------------------------------------
% FICHE DE REMEDIATION — C1 : oubli du +1 dans la mise a jour de borne_inf
%---------------------------------------------------------------

\begin{exercice}{1NSI-ADGK-RE-C1-EX1}{1}{10}
Un élève écrit la recherche dichotomique suivante, avec une erreur dans la mise à jour de
\lstinline{borne_inf} :
\begin{python}
def recherche_dichotomique_bugue(tableau_trie, valeur):
    borne_inf, borne_sup = 0, len(tableau_trie) - 1
    while borne_inf <= borne_sup:
        milieu = (borne_inf + borne_sup) // 2
        if tableau_trie[milieu] == valeur:
            return milieu
        elif tableau_trie[milieu] < valeur:
            borne_inf = milieu  # au lieu de milieu + 1
        else:
            borne_sup = milieu - 1
    return -1
\end{python}
\begin{enumerate}
  \item Sur le tableau \lstinline{[1, 3, 5, 7, 9, 11, 13]}, que se passe-t-il si on cherche la valeur $4$ (absente) avec cette fonction bugée ? (on pourra limiter le nombre d'itérations testées à $20$ pour éviter une exécution trop longue)
  \item Calculer le variant $V=\text{borne\_sup}-\text{borne\_inf}$ à chaque itération. Décroît-il toujours strictement avec cette version bugée ?
  \item Corriger le code pour que la terminaison soit à nouveau garantie.
\end{enumerate}
\end{exercice}

% BEGIN-VERIFY
% def recherche_dichotomique_bugue_limitee(tableau_trie, valeur, max_iter=20):
%     borne_inf, borne_sup = 0, len(tableau_trie) - 1
%     variants = []
%     compteur = 0
%     while borne_inf <= borne_sup:
%         compteur += 1
%         if compteur > max_iter:
%             return "boucle_infinie", variants
%         variants.append(borne_sup - borne_inf)
%         milieu = (borne_inf + borne_sup) // 2
%         if tableau_trie[milieu] == valeur:
%             return milieu, variants
%         elif tableau_trie[milieu] < valeur:
%             borne_inf = milieu
%         else:
%             borne_sup = milieu - 1
%     return -1, variants
% t = [1, 3, 5, 7, 9, 11, 13]
% resultat, variants = recherche_dichotomique_bugue_limitee(t, 4)
% assert resultat == "boucle_infinie"
% assert variants[-1] == variants[-2] == variants[-3]  # le variant ne decroit plus (stagne)
% def recherche_dichotomique(tableau_trie, valeur):
%     borne_inf, borne_sup = 0, len(tableau_trie) - 1
%     while borne_inf <= borne_sup:
%         milieu = (borne_inf + borne_sup) // 2
%         if tableau_trie[milieu] == valeur:
%             return milieu
%         elif tableau_trie[milieu] < valeur:
%             borne_inf = milieu + 1
%         else:
%             borne_sup = milieu - 1
%     return -1
% assert recherche_dichotomique(t, 4) == -1
% END-VERIFY
